\chapter{Applications and Demonstration}
\section{Replay Inference}
The replay application provides a repeatable way to inspect the system. A recorded
video is submitted to the service, processed through the same causal model and
decoder used by the live path, and displayed with the predicted current step.
The interface also shows the procedure list and a human-versus-AI timeline so
that a reviewer can compare recognition quality, transition timing, and
out-of-order behavior on the same recording.
\begin{figure}[H]
\centering
\includegraphics[width=0.94\textwidth]{replay_inference_loaded.png}
\caption{Replay inference interface with a project recording loaded.}
\label{fig:replay}
\end{figure}
\section{Live Inference Path}
The live demonstration receives camera frames over the recording/WebRTC path.
The service derives the model's 10 fps input, maintains a rolling clip buffer,
computes features, updates the causal model, and publishes the current state.
This path exercises transport, frame sampling, feature extraction, temporal
state, decoding, and presentation simultaneously. The measured implementation
runs at approximately 1.14 times playback speed, with an end-to-end
announcement delay of approximately 2.84 seconds.

\section{Live State and Step Progression}

The live service publishes the current committed label, procedure direction,
step progress, and attention status. A prediction that changes for one update
does not immediately change the displayed step; the fixed-lag decoder must
first commit the transition. This is why the interface can update frequently
without making the operator watch a flickering sequence of competing labels.

\section{HoloLens Overlay}
The wearable presentation is implemented as a separate Unity Universal Windows
Platform application for HoloLens. The current project uses Unity 2020.3.28f1
with Microsoft Mixed Reality Toolkit/OpenXR components. The trained model
remains on the inference service; the Unity client polls the service state and
renders the result as a movable, resizable 2D overlay.
\begin{figure}[H]
\centering
\includegraphics[width=0.94\textwidth]{hololens_overlay.jpg}
\caption{HoloLens overlay demonstration showing procedure guidance and current state.}
\label{fig:hololens}
\end{figure}

The HoloLens project is a separate Unity Universal Windows Platform client
using Unity 2020.3.28f1 with Microsoft Mixed Reality Toolkit/OpenXR components.
The trained model remains on the inference service. The Unity client polls the
service state and renders it as a movable, resizable 2D overlay. The current
implementation is screen-space rather than spatially anchored to a physical
part; spatial anchors, part localization, and reference-based geometry checks
are future extensions.

\section{Presentation of the Result}
The user-facing state includes:
\begin{itemize}
    \item the current or most recently committed step;
    \item the ordered procedure list and completion status;
    \item an indication when observed activity is out of order or requires
    attention;
    \item a timeline for comparing human labels with AI commitments.
\end{itemize}

\section{What the Demonstration Proves}

The demonstration proves that the research result has been connected to a
complete application path. It does not prove that every assembly fault can be
detected or that the current screen-space overlay is ready for production use.
The evidence is strongest for the following claims: the system can process
egocentric recordings, maintain a causal running estimate, stabilize its
announcements, display procedure progress, and expose a usable interface for
review or wearable presentation.
